% !PRINTED 268
\weberpage{268}
Demnach wird (11)
\[
(uv)^{\frac{p+1}{2}}\left[\Bigl(\frac{v}{u}\Bigr)^{r\alpha}+\Bigl(\frac{u}{v}\Bigr)^{r\alpha}\right]\left[(uv)^{s\beta}+\frac{\varepsilon^\beta2^{s\beta}}{(uv)^{s\beta}}\right],
\]
und (12)
\[
(uv)^{\frac{p+1}{2}}\left[(uv)^{s\beta}+\frac{\varepsilon^\beta2^{s\beta}}{(uv)^{s\beta}}\right].
\]
Nun gelten die bekannten Formeln, wenn $x, \gamma$ beliebige
Gr\"o\ss en sind und
\[
x+\frac{\gamma}{x}=y
\]
gesetzt wird,
\[
x^2+\frac{\gamma^2}{x^2}=y^2-2\gamma, \qquad x^3+\frac{\gamma^3}{x^3}=y^3-3\gamma y,\ \ldots,
\]
\[
\Bigl(x^n+\frac{\gamma^n}{x^n}\Bigr)\Bigl(x+\frac{\gamma}{x}\Bigr)=x^{n+1}+\frac{\gamma^{n+1}}{x^{n+1}}+\gamma\Bigl(x^{n-1}+\frac{\gamma^{n-1}}{x^{n-1}}\Bigr),
\]
woraus man durch den Schlu\ss\ von $n$ auf $n+1$ erkennt, da\ss
\[
x^n+\frac{\gamma^n}{x^n}
\]
sich f\"ur jedes beliebige $n$ als ganze rationale Funktion $n$ten
Grades von $y$ darstellen l\"a\ss t, die, wenn $\gamma$ eine ganze
Zahl ist, ganzzahlige Koeffizienten hat, deren h\"ochster $=1$ ist.

Die beiden Gr\"o\ss en
\[
\Bigl(\frac{v}{u}\Bigr)^{r\alpha}+\Bigl(\frac{u}{v}\Bigr)^{r\alpha} \qquad \text{und} \qquad (uv)^{s\beta}+\frac{\varepsilon^\beta2^{s\beta}}{(uv)^{s\beta}}
\]
k\"onnen also in dieser Weise als ganze rationale Funktionen von
$A$ und $B$ der Grade $\alpha$ und $\beta$ dargestellt werden, und
wir finden, wenn wir noch das dem Werte $\beta=\dfrac{p-1}{2s}$
entsprechende Glied, das den Koeffizienten $-1$ hat, absondern und
mit $c_{\alpha,\beta}$ ganzzahlige Koeffizienten bezeichnen:
\weq{14}{\frac{\Phi_p(u,v)}{(uv)^{\frac{p+1}{2}}}=A^{\frac{p+1}{2r}}-B^{\frac{p-1}{2s}}+\sum^{\alpha,\beta}c_{\alpha,\beta}A^\alpha B^\beta,}
worin $\alpha$ und $\beta$ an die Grenzen
\[
0\leqq\alpha<\frac{p+1}{2r}, \qquad 0\leqq\beta<\frac{p-1}{2s}
\]
gebunden sind. Dies ist die Form, die wir im vorigen Paragraphen
den Modulargleichungen bis $p=19$ gegeben haben.
